\section*{Hoofdstuk \ref{ch:g7:central} --- Puntspiegeling en parallellogrammen}

\begin{solution}{exo:g7:central:1}
Het beeld onder de halve draai: $M'$ ligt vier vakjes \emph{links} van $O$ en één
vakje \emph{lager} (beide richtingen omgekeerd). Het \omterm{def:g7:central:def}{spiegelbeeld} in de verticale
\omterm{def:g6:lines:objects}{lijn}: vier vakjes links, maar nog steeds één vakje \emph{hoger}. De twee beelden
verschillen: ze zijn verticale \omterm{def:g7:central:def}{spiegelbeelden} van elkaar.
\end{solution}

\begin{solution}{exo:g7:central:2}
$A'B' = AB$ (halve draaien bewaren lengtes) en $(A'B') \parallel (AB)$ (het beeld
van een \omterm{def:g6:lines:objects}{lijn} is een parallelle \omterm{def:g6:lines:objects}{lijn}, \cref{prop:g7:central:preserved}).
\end{solution}

\begin{solution}{exo:g7:central:3}
Cijfers op een display met een \omterm{def:g7:central:center}{symmetriemiddelpunt}: $0$, $2$, $5$, $8$ (en $1$ als je
hem als een streepje tekent). Letters: H, N, S, Z en O hebben er een; A en T niet
(die hebben in plaats daarvan een as).
\end{solution}

\begin{solution}{exo:g7:central:4}
De zijden die tegenover elkaar liggen zijn gelijk: $MN = KL = 7$ cm en
$NK = LM = 4$ cm. De hoeken die tegenover elkaar liggen zijn gelijk:
$\widehat M = \widehat K = 115^\circ$. De diagonalen snijden elkaar in hun
gemeenschappelijke midden, en dat punt is $O$: het midden van $[LN]$ is $O$.
\end{solution}

\begin{solution}{exo:g7:central:5}
Teken een \omterm{def:g6:lines:objects}{lijnstuk} $[AC]$ van $8$ cm, zet zijn midden $O$, teken door $O$ een \omterm{def:g6:lines:objects}{lijn}
die $60^\circ$ met $(AC)$ maakt, en zet $B$ en $D$ erop op $2.5$ cm van $O$, aan elke
kant één. Verbind $A$, $B$, $C$ en $D$: de diagonalen snijden elkaar in hun middens,
dus is $ABCD$ per definitie een \omterm{def:g7:central:parallelogram}{parallellogram}.
\end{solution}

\begin{solution}{exo:g7:central:6}
Opeenvolgende hoeken zijn supplementair: naast $115^\circ$ ligt
$180 - 115 = 65^\circ$. De hoeken die tegenover elkaar liggen zijn gelijk, dus zijn
de vier hoeken $115^\circ$, $65^\circ$, $115^\circ$ en $65^\circ$ (samen
$360^\circ$).
\end{solution}

\begin{solution}{exo:g7:central:7}
Van $A(1,1)$ naar $O(2.5, 2.5)$: $1.5$ naar rechts en $1.5$ naar boven; nog eens
hetzelfde geeft $A'(4, 4)$. Van $B(4,2)$ naar $O$: $1.5$ naar links en $0.5$ naar
boven; nog eens geeft $B'(1, 3)$.
\end{solution}

\begin{solution}{exo:g7:central:8}
Nee. Een \omterm{def:g6:shapes:triangles}{gelijkbenig} trapezium heeft twee gelijke schuine zijden zonder een
\omterm{def:g7:central:parallelogram}{parallellogram} te zijn (de twee parallelle zijden hebben verschillende lengtes). Wat
wel volstaat, is criterium 3 van \cref{thm:g7:central:recognize}: twee zijden die
tegenover elkaar liggen even lang \emph{én} \omterm{def:g6:lines:perp}{parallel}.
\end{solution}

\begin{solution}{exo:g7:central:9}
$O$ is per keuze het midden van $[BC]$, en het is het midden van $[AA']$ door de
constructie van het \omterm{def:g7:central:def}{spiegelbeeld}. De \omterm{def:g4:geometry:polygon}{vierhoek} $ABA'C$ heeft diagonalen $[AA']$ en
$[BC]$ die het midden $O$ delen: criterium 1 (de definitie) maakt er gratis een
\omterm{def:g7:central:parallelogram}{parallellogram} van.
\end{solution}

\begin{solution}{exo:g7:central:10}
``Hoeken tegenover elkaar gelijk $\Rightarrow$ \omterm{def:g7:central:parallelogram}{parallellogram}'': \emph{waar} (dat is
nog een herkenningscriterium; met de hoekensom $360^\circ$ dwingen gelijke paren
tegenover elkaar supplementaire opeenvolgende hoeken af, en dus parallelle zijden
volgens \cref{thm:g7:triangles:equalities}).

``Loodrechte diagonalen $\Rightarrow$ \omterm{def:g6:shapes:quadrilaterals}{ruit}'': \emph{waar} voor een \omterm{def:g7:central:parallelogram}{parallellogram} (de
diagonalen snijden elkaar al in hun middens).

``Eén \omterm{def:g3:shapes:rightangle}{rechte hoek} $\Rightarrow$ rechthoek'': \emph{waar} --- opeenvolgende hoeken zijn
supplementair, dus worden alle vier recht.
\end{solution}

\begin{solution}{exo:g7:central:11}
Wat de \omterm{def:g4:geometry:polygon}{vierhoek} ook is --- zelfs een heel onregelmatige --- de middens $I$, $J$, $K$
en $L$ vormen altijd een \emph{\omterm{def:g7:central:parallelogram}{parallellogram}}. Het bewijs met de
middenparallelstelling staat in \cref{ch:g8:midpoints}
(\cref{exo:g8:midpoints:9}).
\end{solution}

\begin{solution}{pb:g7:central:1}
\textbf{1.} $M(3, 1) \to (-3, -1)$; $N(-2, 4) \to (2, -4)$;
$P(0, -5) \to (0, 5)$. Regel: de halve draai om de oorsprong stuurt $(x, y)$ naar
$(-x, -y)$ --- beide coördinaten wisselen van teken.

\textbf{2.} Vanaf $M(5, 3)$: drie vakjes rechts van $C$ wordt drie vakjes links, twee
omhoog wordt twee omlaag: het beeld is $(-1, -1)$. Controle met het recept: twee keer
het middelpunt min het punt, $(2 \times 2 - 5,\ 2 \times 1 - 3) = (-1, -1)$.

\textbf{3.} De halve draai om het middelpunt van de kaart wisselt de symbolen in
paren om. Bij een \omterm{def:g3:numbers:evenodd}{oneven} aantal symbolen heeft één symbool geen partner: het moet zijn
eigen beeld zijn, en het enige punt dat zijn eigen beeld is, is het middelpunt zelf.
Het middelste symbool van de 3, de 5, de 7\,\dots\ ligt dus precies in het midden
(kijk naar een echte kaart: het klopt).

\textbf{4.} Om $O_1(2,0)$: $A(1,1) \to (3,-1)$; om $O_2(5,0)$:
$(3,-1) \to (7,1)$. Begin $A(1,1)$, einde $(7,1)$: $6$ naar rechts geschoven, niet
omgeklapt. Net zo voor $B(0,3) \to (4,-3) \to (6,3)$: $6$ naar rechts. De twee halve
draaien komen neer op een \emph{verschuiving}.

\textbf{5.} De verschuiving is $6$, en $O_1 O_2 = 3$: de verschuiving legt
\emph{twee keer} de afstand tussen de middelpunten af (in de richting van $O_1$ naar
$O_2$) --- net zoals twee parallelle spiegels in \cref{pb:g6:symmetry:1} over twee
keer hun tussenafstand verschoven.

\textbf{6.} De halve draai om $O$ beeldt het \omterm{def:g7:central:parallelogram}{parallellogram} op zichzelf af (dat is
wat ``\omterm{def:g7:central:center}{symmetriemiddelpunt}'' er voor betekent,
\cref{thm:g7:central:properties}). Een \omterm{def:g6:lines:objects}{lijn} door $O$ wordt ook op zichzelf
afgebeeld, dus wisselen de twee punten waar ze de rand snijdt met elkaar om: ze zijn
\omterm{def:g7:central:def}{spiegelbeelden} om $O$. De twee stukken van het \omterm{def:g7:central:parallelogram}{parallellogram} aan weerszijden van de
\omterm{def:g6:lines:objects}{lijn} wisselen ook om, en zijn dus precies elkaars kopie
(\cref{prop:g7:central:preserved}) --- gelijke \omterm{def:g6:measure:area}{oppervlaktes}.

\textbf{7.} Snijd langs de \omterm{def:g6:lines:objects}{lijn} door de \emph{twee middelpunten}: het middelpunt van
de rechthoekige cake en het middelpunt van het gat. Volgens vraag 6, toegepast op de
cakerechthoek, halveert de snede de cake; toegepast op de gatrechthoek (de \omterm{def:g6:lines:objects}{lijn} gaat
ook door haar middelpunt) halveert ze het gat. Elk deel krijgt de \omterm{def:g3:division:half}{helft} van de cake
min de \omterm{def:g3:division:half}{helft} van het gat: volmaakt eerlijk.

\textbf{8.} $D = (0 + 8 - 6,\ 0 + 5 - 1) = (2, 4)$ (de diagonalen delen hun midden,
dus is $D$ het beeld van $B$ onder de halve draai om het middelpunt). Middelpunt: het
midden van $[AC]$, dus $O(4,\ 2.5)$.

\textbf{9.} In het \omterm{def:g7:central:parallelogram}{parallellogram} $ABA'C$ zijn de zijden die tegenover elkaar liggen
gelijk: $BA' = AC$. De driehoeksongelijkheid in $ABA'$ geeft
$AA' < AB + BA' = AB + AC$. Omdat $O$ het midden van $[AA']$ is, is $AA' = 2\,AO$,
dus $2\,AO < AB + AC$, en dus $AO < \frac{AB + AC}{2}$.

\textbf{10.} Bovengrens: $AO < \frac{5 + 7}{2} = 6$ cm. Ondergrens: in $ABA'$ is
$AA' > BA' - AB = 7 - 5 = 2$, dus $AO > 1$ cm. De zwaartelijn uit $A$ ligt strikt
tussen $1$ en $6$ cm.

\textbf{11.} \omterm{def:g6:lines:objects}{Lijnstuk}: ja, zijn midden. \omterm{def:g6:lines:objects}{Lijn}: ja --- elk van haar punten is een
middelpunt. \omterm{def:g6:shapes:triangles}{Gelijkzijdige driehoek}: geen. Vierkant: ja, het snijpunt van de
diagonalen. Cirkel: ja, zijn middelpunt. S en Z: ja, hun middelpunt (draai de
bladzijde om: ze lezen hetzelfde).

\textbf{12.} De halve draai zou de drie \omterm{def:g5:solids:def}{hoekpunten} onder elkaar in paren zetten;
omdat drie \omterm{def:g3:numbers:evenodd}{oneven} is, zou één \omterm{def:g5:solids:def}{hoekpunt} $A$ zijn eigen beeld zijn, en dan moet $A$ het
middelpunt van de halve draai zijn. De twee andere \omterm{def:g5:solids:def}{hoekpunten} $B$ en $C$ zouden dan
\omterm{def:g7:central:def}{spiegelbeelden} om $A$ zijn --- dus zou $A$ het midden van $[BC]$ zijn, en zouden de
drie \omterm{def:g5:solids:def}{hoekpunten} op één \omterm{def:g6:lines:objects}{lijn} liggen. Een driehoek heeft geen drie \omterm{def:g5:solids:def}{hoekpunten} op één
\omterm{def:g6:lines:objects}{lijn}: tegenspraak. Er bestaat dus geen middelpunt.

\textbf{13.} Hetzelfde pariteitsargument: een halve draai die de \omterm{def:g4:geometry:polygon}{veelhoek} bewaart,
zet haar \omterm{def:g5:solids:def}{hoekpunten} in paren; een \omterm{def:g3:numbers:evenodd}{oneven} aantal \omterm{def:g5:solids:def}{hoekpunten} dwingt een vast \omterm{def:g5:solids:def}{hoekpunt}
af, dat het middelpunt moet zijn en tevens het midden van het \omterm{def:g6:lines:objects}{lijnstuk} tussen twee
andere \omterm{def:g5:solids:def}{hoekpunten} --- drie \omterm{def:g5:solids:def}{hoekpunten} op één \omterm{def:g6:lines:objects}{lijn}, onmogelijk in een \omterm{def:g4:geometry:polygon}{veelhoek}\,\dots\
dus heeft geen \omterm{def:g4:geometry:polygon}{veelhoek} met een \omterm{def:g3:numbers:evenodd}{oneven} aantal \omterm{def:g5:solids:def}{hoekpunten} een \omterm{def:g7:central:center}{symmetriemiddelpunt}.

\textbf{14.} De halve draai om $O_1$ en daarna die om $O_2$ beelden de figuur beide op
zichzelf af --- dus doet hun samenstelling dat ook. Volgens vraag 5 is die
samenstelling een verschuiving over twee keer de afstand $O_1 O_2$, en die is niet \omterm{def:g1:counting:zero}{nul}
omdat de middelpunten verschillen. Een figuur die op een blad past, kan niet gelijk
zijn aan zichzelf over een vaste afstand verschoven (schuif haar vaak genoeg en ze
verlaat het blad volledig). Een begrensde figuur heeft dus hoogstens één
\omterm{def:g7:central:center}{symmetriemiddelpunt}.

\textbf{15.} Op de eindeloze strook beeldt een halve draai om het punt halverwege
tussen een linkervoetstap en de volgende rechtervoetstap het patroon op zichzelf af;
het punt één volle stap verder doet dat ook: twee verschillende middelpunten. De twee
halve draaien samenstellen geeft de verschuiving over twee keer die halve stap ---
precies de verschuiving van één stap die de oneindige fries zichtbaar op zichzelf
afbeeldt, zoals vraag 14 voorspelt. Onbegrensde patronen leven volgens andere regels:
dat is de wiskunde van het behangpapier.
\end{solution}
